% ---- score2 standard preamble (from: python3 score2_tools.py preamble) ----
\documentclass[10pt]{extarticle}
\usepackage[a4paper,margin=1.5cm]{geometry}
\usepackage[T1]{fontenc}
\usepackage{amsmath,amssymb}
\usepackage{graphicx}
\usepackage[export]{adjustbox}
\usepackage{multicol}
\usepackage{enumitem}
\usepackage{booktabs,array,tabularx}
\setlength{\columnsep}{0.7cm}
\setlength{\parindent}{0pt}
\setlength{\parskip}{3pt}
\pagestyle{empty}
\makeatletter
% figure registry (figdefs.tex is generated by the extract step)
\newcommand{\qfigdef}[4]{\@namedef{qf@file@#1}{#2}\@namedef{qf@w@#1}{#3}\@namedef{qf@wide@#1}{#4}}
\newcommand{\qf@check}[1]{\@ifundefined{qf@file@#1}{\PackageError{score2}{Unknown figure id #1}{}}{}}
\newcommand{\qf@img}[4][]{\includegraphics[width=#2\dimexpr\@nameuse{qf@w@#4}pt\relax,max width=#3,max height=0.45\textheight,#1]{\@nameuse{qf@file@#4}}}
% Full-width questions: finalize wraps every question that contains a wide figure in \qfullbegin...\qfullend,
% so the WHOLE question (stem, figures, options) is set across the page in reading order, then the two columns
% resume. The question's head - its number, stem and first figure (\qfullsplit marks the end of the head) -
% is kept on one page, so a number is never left at a page foot with its figure overleaf; the rest may break
% across pages normally, so no big gaps are left. Without a \qfullsplit, the whole question is kept together
% if it is under 45% of a page. Never write these yourself.
\newif\ifqf@full
\newif\ifqf@boxed
\newsavebox{\qf@box}
\newcommand{\qf@open}{\global\qf@fulltrue\par\global\qf@boxedtrue
  \setbox\qf@box\vbox\bgroup\hsize\textwidth\linewidth\textwidth}
\newcommand{\qf@emit}[1]{\par\egroup\global\qf@boxedfalse
  \ifdim\dimexpr\ht\qf@box+\dp\qf@box\relax<#1\textheight\noindent\usebox{\qf@box}\par\else\unvbox\qf@box\fi}
\newcommand{\qfullsplit}{\ifqf@boxed\qf@emit{0.85}\fi}
\newcommand{\qf@close}{\ifqf@boxed\qf@emit{0.45}\fi\global\qf@fullfalse}
\newcommand{\qfullbegin}{\end{multicols}\qf@open}
\newcommand{\qfullend}{\qf@close\begin{multicols}{2}}
\newcommand{\qfullstart}{\qf@open}
\newcommand{\qfulldone}{\qf@close}
% \qfig[scale]{id}: block figure, centred, directly where it is written (size set by the tool).
% Fallback if a wide figure is not inside a full-width question: break out of the columns in place - no floats.
\newcommand{\qfig}[2][1]{\qf@check{#2}%
  \ifqf@full\qf@blk{#1}{#2}\else
  \ifnum\@nameuse{qf@wide@#2}=1 \expandafter\@firstoftwo\else\expandafter\@secondoftwo\fi
  {\end{multicols}\begin{center}\qf@img{#1}{\textwidth}{#2}\end{center}\begin{multicols}{2}}%
  {\qf@blk{#1}{#2}}\fi}
\newcommand{\qf@blk}[2]{\par\smallskip{\centering\qf@img{#1}{\linewidth}{#2}\par}\smallskip}
% \qfigi[scale]{id}: inline figure (inside option lists / tables / minipages); never breaks columns.
\newcommand{\qfigi}[2][1]{\qf@check{#2}\qf@img[valign=c]{#1}{\linewidth}{#2}}
\makeatother
% structure macros - lint reads these to check completeness and figure placement
\newcommand{\papersec}[1]{\par\bigskip{\centering\Large\bfseries PAPER-#1\par}\medskip}
\newcommand{\partsec}[2]{\par\medskip\noindent{\large\bfseries PART-#1: #2}\par\smallskip}
\newcommand{\sect}[1]{\par\smallskip\noindent\textbf{#1}\par}
\newcommand{\qn}[1]{\par\smallskip\noindent\textbf{#1.}~}
\newcommand{\paraq}[2]{\par\smallskip\noindent\textbf{Paragraph for Questions #1 and #2}\par}
\newenvironment{qoptions}{\begin{itemize}[leftmargin=1.8em,itemsep=1pt,topsep=2pt,label={}]}{\end{itemize}}
\newcommand{\opt}[2]{\item[(#1)] #2}
% figure registry (generated by score2_tools.py finalize)
\qfigdef{p1_phy_s1_q02_1}{Score-II - Test 35 (04 Feb 2025)-img/p1_phy_s1_q02_1.png}{112.6}{0}
\qfigdef{p1_phy_s1_q03_1}{Score-II - Test 35 (04 Feb 2025)-img/p1_phy_s1_q03_1.png}{72.0}{0}
\qfigdef{p1_phy_s2_q06_1}{Score-II - Test 35 (04 Feb 2025)-img/p1_phy_s2_q06_1.png}{75.7}{0}
\qfigdef{p1_phy_s2_q07_1}{Score-II - Test 35 (04 Feb 2025)-img/p1_phy_s2_q07_1.png}{135.9}{0}
\qfigdef{p1_phy_s2_q08_1}{Score-II - Test 35 (04 Feb 2025)-img/p1_phy_s2_q08_1.png}{128.1}{0}
\qfigdef{p1_phy_s2_q09_1}{Score-II - Test 35 (04 Feb 2025)-img/p1_phy_s2_q09_1.png}{143.7}{0}
\qfigdef{p1_phy_s3_q01_1}{Score-II - Test 35 (04 Feb 2025)-img/p1_phy_s3_q01_1.png}{49.8}{0}
\qfigdef{p1_phy_s3_q02_1}{Score-II - Test 35 (04 Feb 2025)-img/p1_phy_s3_q02_1.png}{81.5}{0}
\qfigdef{p1_phy_s3_q03_1}{Score-II - Test 35 (04 Feb 2025)-img/p1_phy_s3_q03_1.png}{53.3}{0}
\qfigdef{p1_phy_s3_q04_1}{Score-II - Test 35 (04 Feb 2025)-img/p1_phy_s3_q04_1.png}{112.2}{0}
\qfigdef{p1_phy_s3_q05_1}{Score-II - Test 35 (04 Feb 2025)-img/p1_phy_s3_q05_1.png}{135.5}{0}
\qfigdef{p1_phy_s3_q06_1}{Score-II - Test 35 (04 Feb 2025)-img/p1_phy_s3_q06_1.png}{159.2}{0}
\qfigdef{p1_phy_s3_q07_1}{Score-II - Test 35 (04 Feb 2025)-img/p1_phy_s3_q07_1.png}{50.3}{0}
\qfigdef{p1_phy_s3_q08_1}{Score-II - Test 35 (04 Feb 2025)-img/p1_phy_s3_q08_1.png}{105.2}{0}
\qfigdef{p1_chem_s2_q06_1}{Score-II - Test 35 (04 Feb 2025)-img/p1_chem_s2_q06_1.png}{135.4}{0}
\qfigdef{p1_chem_s2_q08_1}{Score-II - Test 35 (04 Feb 2025)-img/p1_chem_s2_q08_1.png}{156.7}{0}
\qfigdef{p1_chem_s2_q10_1}{Score-II - Test 35 (04 Feb 2025)-img/p1_chem_s2_q10_1.png}{167.5}{0}
\qfigdef{p1_chem_s3_q01_1}{Score-II - Test 35 (04 Feb 2025)-img/p1_chem_s3_q01_1.png}{232.1}{0}
\qfigdef{p1_chem_s3_q06_1}{Score-II - Test 35 (04 Feb 2025)-img/p1_chem_s3_q06_1.png}{205.2}{0}
\qfigdef{p1_chem_s3_q08_1}{Score-II - Test 35 (04 Feb 2025)-img/p1_chem_s3_q08_1.png}{221.0}{0}
\raggedcolumns
\makeatletter\@beginparpenalty=10000\makeatother
\begin{document}
\begin{center}{\Large\bfseries JEE (Main+Advanced) Score-II --- Test 35 (04 Feb 2025)}\end{center}
\begin{multicols}{2}

\partsec{1}{PHYSICS}
\sect{Section-I(i)}

\qn{1} In Young's double slit experiment, the distance d between the slits $S_1$ and $S_2$ is 1.0 mm. What should be the width of each slit so as to obtain 5 maxima of the two slit interference pattern within the central maximum of the single slit diffraction pattern?
\begin{qoptions}
\opt{A}{0.2 mm}
\opt{B}{0.4 mm}
\opt{C}{12 cm}
\opt{D}{0.1 mm}
\end{qoptions}

\qn{2} Figure shows an amperian path ABCDA. Part ABC is in vertical plane PSTU while part CDA is in horizontal plane PQRS. Direction of circulation along the path is shown by an arrow near point B and at D. $\oint\vec B\cdot d\vec l$ for this path according to Ampere's law will be :-
\qfig{p1_phy_s1_q02_1}
\begin{qoptions}
\opt{A}{$(i_1-i_2+i_3)\mu_0$}
\opt{B}{$(-i_1+i_2)\mu_0$}
\opt{C}{$i_3\mu_0$}
\opt{D}{$(i_1+i_2)\mu_0$}
\end{qoptions}

\qn{3} A spherical cave of radius R/2 was carved out from a uniform sphere of radius R and original mass M. The center of the cave is at R/2 from the center of the large sphere. Point P is at a distance 2R from the center of the large sphere and on the joint line of the two centers. The magnitude of gravitational field strength g at point P is \underline{\hspace{1cm}}.
\qfig{p1_phy_s1_q03_1}
\begin{qoptions}
\opt{A}{$g=\dfrac{GM}{4R^2}$}
\opt{B}{$g=\dfrac{5GM}{8R^2}$}
\opt{C}{$g=\dfrac{3GM}{16R^2}$}
\opt{D}{$g=\dfrac{7GM}{36R^2}$}
\end{qoptions}

\qn{4} In a circuit for finding the resistance of a galvanometer by half deflection method, a 6 V battery and a high resistance of 11 k$\Omega$ are used. The figure of merit of the galvanometer is $60\,\mu$A/division. In the absence of shunt resistance, the galvanometer produces a deflection of $\theta=9$ divisions when current flows in the circuit. The value of the shunt resistance that can cause the deflection of $\theta/2$, is closest to :-
\begin{qoptions}
\opt{A}{55 $\Omega$}
\opt{B}{110 $\Omega$}
\opt{C}{220 $\Omega$}
\opt{D}{550 $\Omega$}
\end{qoptions}

\sect{Section-I(ii)}

\qn{5} Three nonconducting infinite planar sheets are parallel to the y-z plane. Each sheet has a uniform surface charge density. The first sheet, with a negative surface charge density $-\sigma$, passes through the x-axis at x = 1 m. The second sheet has an unknown surface charge density and passes through the x-axis at x = 2m. The third sheet has a negative surface charge density $-3\sigma$ and passes through the x-axis at x = 4 m. The net electric field due to the sheets is zero at x = 1.5 m.
\begin{qoptions}
\opt{A}{The surface charge density on the second sheet is $+2\sigma$.}
\opt{B}{The electric field at x = $-2$ m is $\dfrac{\sigma}{\epsilon_0}\hat i$}
\opt{C}{The electric field at x = 3 m is $\dfrac{\sigma}{\epsilon_0}\hat i$}
\opt{D}{The electric field at x = 6m is $\dfrac{-\sigma}{\epsilon_0}\hat i$}
\end{qoptions}

\qn{6} A time varying magnetic field exists in a cylindrical region of radius a. A semi circular conductor of negligible resistance with a small voltmeter $V_2$ of resistance 2R is placed coaxially with cylindrical region as shown in diagram. A small voltmeter $V_1$ having resistance R is connected across the two ends of semicircular conductor as shown in diagram. Magnetic field is varying at a rate of $\dfrac{dB}{dt}$. Choose correct option(s).
\qfig{p1_phy_s2_q06_1}
\begin{qoptions}
\opt{A}{Current passing through voltmeter $V_1$ is $\dfrac{\pi a^2}{6R}\cdot\dfrac{dB}{dt}$}
\opt{B}{Reading of voltmeter $V_1$ is zero because it is connected to a conductor which along diameter of cylinderical region.}
\opt{C}{Reading of voltmeter is $V_1$ is $\dfrac{\pi a^2}{6}\cdot\dfrac{dB}{dt}$}
\opt{D}{Reading of voltmeter $V_2$ is $\dfrac{\pi a^2}{3}\cdot\dfrac{dB}{dt}$}
\end{qoptions}

\qn{7} A light ray coming from medium of refractive index $\mu_1$ along the line y = 8 is incident on the interface of a medium of refractive index $\mu_2$ whose curve is given by $y^2=16x$. Then just after first refraction :
\qfig{p1_phy_s2_q07_1}
\begin{qoptions}
\opt{A}{If $\mu_1=1$ and $\mu_2=\sqrt2$ then angle of deviation is $\dfrac\pi{12}$.}
\opt{B}{If $\mu_1=1$ and $\mu_2=\sqrt2$ then angle of deviation is $\dfrac{5\pi}{12}$.}
\opt{C}{If $\mu_1=\sqrt2$ and $\mu_2=1$ then angle of deviation is $\dfrac\pi4$.}
\opt{D}{If $\mu_1=\sqrt2$ and $\mu_2=1$ then angle of deviation is $\dfrac\pi2$.}
\end{qoptions}

\qn{8} A car is moving with a velocity $\dfrac{3v}{10}$ towards a medium MNRP as shown, where `v' is the speed of sound in air. Two men A and B are moving towards each other in the medium with the velocities $\dfrac v6$ and $\dfrac v9$ respectively. The speed of sound in medium is 7v and medium is kept at rest. The car is emitting a sound of frequency f. Then choose the correct option(s).
\qfig{p1_phy_s2_q08_1}
\begin{qoptions}
\opt{A}{The wavelength received by the man A is $\dfrac{49v}{10f}$.}
\opt{B}{The wavelength received by the man A is different form the wavelength received by man B.}
\opt{C}{The ratio of frequency observed by the man A to that observed by the man B is $\dfrac{123}{128}$}
\opt{D}{The frequency observed by the man B is $\left(\dfrac{640}{441}\right)$f.}
\end{qoptions}

\qn{9} A uniform rod of mass m and length $\ell$ is hinged at point P such that the rod can rotate in a vertical plane. One end B of the rod is connected through the massless string which passes over the pulley of mass m and radius R and the other end of the string is connected with the block of mass m as shown in the figure. (There is no slipping between the string and the pulley) Initially rod is horizontal. At t = 0, the system is released, then which of the following is/are correct.
\qfig{p1_phy_s2_q09_1}
\begin{qoptions}
\opt{A}{Acceleration of the block is $\dfrac{36}{95}$g}
\opt{B}{Acceleration of the centre of mass of the rod is $\dfrac{12}{95}$g}
\opt{C}{Tension in the string $S_1$ is $\dfrac{41}{95}$mg}
\opt{D}{Tension in the string $S_2$ is $\dfrac{39}{95}$mg}
\end{qoptions}

\qn{10} A small metal plate of area $2\ \text{cm}^2$ is placed at a distance of 2 m from a monochromatic light source of $\lambda=1000$ \AA{} and power rating of 100 W and light energy conversion efficiency of 10\%, emitting light uniformly in all directions. Light falls normally on the plate (work function, $\phi=10$eV). Assuming one out of $10^6$ photons is able to eject a photoelectron, identify correct statements
\begin{qoptions}
\opt{A}{KE max of photoelectron emitted is 2.4 eV}
\opt{B}{Saturation current is $3.2\times10^{-12}$ A}
\opt{C}{Saturation current is $4.5\times10^{-12}$ A}
\opt{D}{KE max of photoelectron emitted is 3.2 eV}
\end{qoptions}

\sect{Section-III}

\qn{1} A solid object having uniform density $(\rho)$ and varying cross-sectional area (which is circular) is placed on smooth horizontal surface as shown. If the stress throughout its height remains same and having value equal to c, the value of $R_2=R_1e^{\frac{\rho gH}{nc}}$. The value of n is :-
\qfig{p1_phy_s3_q01_1}

\qn{2} The magnetic field existing in a region is given by $\vec B=0.2\,(1+2x)\,\hat k$T. A square loop of edge 50 cm carrying 0.5 A current is placed in x-y plane with its edges parallel to the x-y axes, as shown in figure. The magnitude of the net magnetic force experienced by the loop is $\alpha$ mN. Find the value of $\alpha/10$, where $\alpha$ is a least integer.
\qfig{p1_phy_s3_q02_1}

\qn{3} The shaded area in the figure is the side view of a disk shaped magnetic field region of radius R and thickness d. A parallel electron beam being accelerated by a voltage V is normally incident onto the region. If the spatial distribution of the magnetic field in the region such that the beam is focused to a point along the disk central axis at a distance L ($\gg$ R) from the disk is $B=\dfrac{y^\eta}{d^\alpha L^\beta}\sqrt{\dfrac{\lambda mV}{e}}$, find the value of $\dfrac{20\alpha}{\lambda+\beta+\eta}$. Given m and e are the electron's mass and charge and y is the distance of incident electron from axis of disc.
\qfig{p1_phy_s3_q03_1}

\qn{4} A magnetized needle of length 10 cm is freely suspended by its south pole from a hook as shown in figure. Mass of needle is 2 milligram. In equilibrium, the needle makes an angle of $60^\circ$ from vertical. The magnetic pole strength is concentrated at the ends of the needle. The magnetic pole strength of needle (in A-m) is $\beta$. Find the value of $10\beta$. (Given : number of turns in coil = 1,000, radius of coil = 10 cm, current in coil $=\dfrac{100}{\pi}$ A and end of needle is lying on the axis of coil at a distance of 2 m from the centre of coil.)
\qfig{p1_phy_s3_q04_1}

\qn{5} The given system is in equilibrium. The velocity of wave in strings A \& B is respectively x and y (in m/s). Given the answer $\dfrac{x+y}{15}$.
\qfig{p1_phy_s3_q05_1}

\qn{6} In a particular set up, interference is observed on an screen as shown. Wavelength of light emitted by source is 4000 \AA. Out of two coherent sources, one is obtained by reflection with a plane mirror of length 2m. The number of points on the screen where intensity is minimum, is :
\qfig{p1_phy_s3_q06_1}

\qn{7} A small current carrying loop having current I is placed in the plane of paper as shown. Another semicircular loop having current $I_0$ is placed concentrically in the same plane as that of small loop, the radius of semi-circular loop is R(R$\gg$a). Find the force applied by the smaller ring on bigger ring in newton.\\
$\left(\text{Given R = 1m, I = }I_0=\dfrac{40}{\sqrt{\mu_0}}\text{ A, a = 0.1 m}\right)$
\qfig{p1_phy_s3_q07_1}

\qn{8} The figure shows a 0.9 m wide uniform door hinged at the edge being subjected to a force of magnitude 12 N as shown. About the given axis, the moment of inertia of the door is 1.2 kg m$^2$. Neglecting gravity, what is the force (in N) exerted by the hinge as soon as the force is applied ?
\qfig{p1_phy_s3_q08_1}

\partsec{2}{CHEMISTRY}
\sect{Section-I(i)}

\qn{1} Concentration of H$_2$SO$_4$ and Na$_2$SO$_4$ in a solution is 1 M and $1.8\times10^{-2}$ M, respectively. Molar solubility of PbSO$_4$ in the same solution is :\\
\relax[Given: Solubility product of PbSO$_4$(K$_{sp}$) = $1.6\times10^{-8}$. For H$_2$SO$_4$, Ka$_1$ is very large and Ka$_2$ = $1.2\times10^{-4}$]
\begin{qoptions}
\opt{A}{$\dfrac89\times10^{-6}$ M}
\opt{B}{$\sqrt{1.6}\times10^{-4}$ M}
\opt{C}{$1.33\times10^{-4}$ M}
\opt{D}{$1.6\times10^{-6}$ M}
\end{qoptions}

\qn{2} An element ``X'' (atomicity = 2) forms oxide X$_M$O$_N$ \& X$_P$O$_R$ (in significant amount) with oxygen during burning of fossil fuel at high temperature in an automobile engine. The oxide X$_P$O$_R$ is a lung irritant \& usually causes respiratory problems in children.\\
Which of the following is correct regarding X$_M$O$_N$ \& X$_P$O$_R$?
\begin{qoptions}
\opt{A}{Both X$_M$O$_N$ \& X$_P$O$_R$ are diamagnetic.}
\opt{B}{Both X$_M$O$_N$ \& X$_P$O$_R$ have tendency to dimerize.}
\opt{C}{Mixture of X$_M$O$_N$ \& X$_P$O$_R$ when cooled yields blue color liquid.}
\opt{D}{Both X$_M$O$_N$ \& X$_P$O$_R$ follow octet rule.}
\end{qoptions}

\qn{3} The complex compound that shows paramagnetic and is a low spin complex among the following.
\begin{qoptions}
\opt{A}{K$_3$[Co(NO$_2$)$_6$]}
\opt{B}{(NH$_4$)$_2$[PtCl$_6$]}
\opt{C}{K$_2$[Co(CN)$_6$]}
\opt{D}{K$_2$[NiF$_6$]}
\end{qoptions}

\qn{4} Isopentane $\xrightarrow[h\nu]{\text{Br}_2}$ monobromo products.\\
Select correct statement about mono bromo product.
\begin{qoptions}
\opt{A}{Total monobromo products are 6 in which molecules of major products are optically active.}
\opt{B}{Total monobromo products are 6 in which major product is optically inactive.}
\opt{C}{Total monobromo products are 4 in which molecules of major products are optically active.}
\opt{D}{Total monobromo products are 4 in which major product is optically inactive.}
\end{qoptions}

\sect{Section-I(ii)}

\qn{5} Select correct statement(s) :
\begin{qoptions}
\opt{A}{Orbital angular momentum depends on two quantum number is $\ell$ and $m_\ell$.}
\opt{B}{Total energy does not very in 1s orbital of H-atom with distance from nucleus.}
\opt{C}{Whenever a radiation intacts with matter it displays wave nature while when it propagates it exhibit particle nature.}
\opt{D}{Boundary surface diagram of p atomic orbitals is dumbell but $|\psi^2|$ is not same at different points on boundary surface.}
\end{qoptions}

\qn{6}
\qfig{p1_chem_s2_q06_1}
X = Mole fraction of liquid A in liquid solution.\\
Y = Mole fraction of vapour A in vapour mixture.\\
Length of line segments KL and LM are in the ratio of 2 : 3.\\
\relax[Given : $P_T$ = 480 mm Hg, $P^\circ_A$ = 300 mm of Hg, $P^\circ_B$ = 600 mm of Hg]
\begin{qoptions}
\opt{A}{X is 0.4}
\opt{B}{Y is 0.25}
\opt{C}{P is 428.57 mmHg.}
\opt{D}{Moles of A in vapour phase at 480 mmHg is 0.4, if total moles of A and B are 1 and 3 respectively.}
\end{qoptions}

\qn{7} The possible type of ions that can be produced by oxides of Vanadium on treatment with acidic or alkaline solutions is/are.
\begin{qoptions}
\opt{A}{VO$^+$}
\opt{B}{VO$^{2+}$}
\opt{C}{VO$_2^+$}
\opt{D}{VO$_4^{3-}$}
\end{qoptions}

\qn{8}
\qfig{p1_chem_s2_q08_1}
The above apparatus shows a sample kept in presence of an electromagnet against some amount of weight on the right side.\\
Assume weight of all samples used in different experiment is same then which of the following sample(s) will require significantly more weight to be kept on right side to balance as compared to the weight required to balance the sample of H$_2$O($\ell$). (The sample can be in either solid or liquid or gas)
\begin{qoptions}
\opt{A}{C$_6$H$_6$($\ell$)}
\opt{B}{K$_4$[Fe(CN)$_6$](s)}
\opt{C}{K$_3$[Fe(CN)$_6$](s)}
\opt{D}{ClO$_2$(g)}
\end{qoptions}

\qn{9} Sucralose is artificial sweetner, the preparation of sucralose involves a multi-step chemical reaction where sucrose is selectively chlorinated by substituting three specific hydroxyl groups by 3 atoms of Cl.
\begin{center}
Sucrose $\rightarrow$ $\underset{\text{Trichloro derivative of sucrose}}{\text{Sucralose}}$
\end{center}
Select the correct statement(s) about sucralose.
\begin{qoptions}
\opt{A}{It is 600 times sweeter than cane sugar.}
\opt{B}{It has 9 chiral carbons.}
\opt{C}{By weight 40\% Cl is present in sucralose.}
\opt{D}{Sucrose and sucralose both do not give positive test with Tollen's reagent.}
\end{qoptions}

\qn{10} An organic acid A (C$_{11}$H$_{12}$O$_2$) can easily be oxidized by H$^{\oplus}$, KMnO$_4$ to a dibasic acid which reacts with ethyleneglycol to produce a polymer dacron. Upon ozonolysis, A gives an aliphatic ketone as one of the product.\\
A gives product B with following reagents.
\qfig{p1_chem_s2_q10_1}
Correct statement(s) about A is/are :
\begin{qoptions}
\opt{A}{It gives positive bromine water test.}
\opt{B}{It gives effervescence with NaHCO$_3$.}
\opt{C}{It give violet coloured complex with neutral FeCl$_3$.}
\opt{D}{It gives positive test with 2,4-Dinitrophenyl hydrazine (2,4-DNP)}
\end{qoptions}

\sect{Section-III}

\qn{1} The most popular electrochemical cell, Daniel cell was originally developed by the English chemist John F. Daniel. The general assembly of the cell is as given below:
\qfig{p1_chem_s3_q01_1}
An undergraduate student made up a Daniel cell using 100 cm$^3$ each of 0.100 M CuSO$_4$ and 0.100 M ZnSO$_4$ solution respectively. The two compartments are connected by suitable salt bridge.\\
\relax[Given : $E^0_{\text{Cu}^{2+}/\text{Cu}}=0.34$ V, $E^0_{\text{Zn}^{2+}/\text{Zn}}=-0.76$ V, $\dfrac{2.303RT}{F}=0.06$, $\log2=0.3$]\\
A lab mate of the student asked her for some solid CuCl$_2$. While she was lifting the bottle from a shelf, the lid of the bottle slipped and some amount of CuCl$_2$ fell in the CuSO$_4$ compartment without changing the volume. She measured the emf of the cell again and found that it had increased by 9 mV. She used this data to calculate the amount of CuCl$_2$ that had split in the compartment. Calculate the mass of CuCl$_2$ that had split into the Daniel cell ?\\
If the answer is `x' gm then the value of $\dfrac{x}{0.45}$ is \underline{\hspace{1cm}}.\\
(Molecular mass of CuCl$_2$ = 135 g)

\qn{2} A vertical cylinder closed from both ends is equipped with an easily moving piston dividing the volume into two parts, each containing one mole of air. In equilibrium at 320 K, the volume of the upper part is 4.0 times greater than that of the lower part. At what temperature (in kelvin) will the ratio of volume of upper part to lower part will be 3 ?\\
If the answer is T, then the value of $\left(\dfrac T{50}\right)$ is.

\qn{3} An athlete in a gymnasium room lifts a 50 kg mass through a vertical distance of 2.0 m; g = 9.8 ms$^{-2}$. The mass is allowed to fall through 2.0 m distance while coupled to an electrical generator. The electrical generator produces an equal amount of electrical work. This electrical work is used to produce aluminum by Hall's process involving the change,
\begin{center}
Al$_2$O$_3$ (molten) + 3C(s) $\rightarrow$ 2Al (l) + 3CO(g); $\Delta G^\circ$ = 588 kJ/mol.
\end{center}
How many times (in Hundred) must the athlete lift the 50 kg mass and drop to couple the generator to produce sufficient Gibbs free energy to produce 27g Al ?

\qn{4} Three moles of copper(II) sulphide is treated with excess dilute HNO$_3$ to yield a precipitate ``X'' and gas ``Y''. ``X'' is then separately treated with excess dilute HNO$_3$ again along with heating to yield gas ``Y'' again. The total moles of gas ``Y'' produced will be.

\qn{5} Total number of degenerate orbitals in $2^{\text{nd}}$ excited state of Be$^{3+}$ will be \underline{\hspace{1cm}}.\\
(Assume absence of external electric and magnetic filed)

\qn{6} How many total statements are correct about DNA.\\
(i) DNA has specific sequence of nucleotides.\\
(ii) DNA synthesis takes place from $5'$ end to $3'$ end.\\
(iii) DNA double helix model is the secondary structure of DNA\\
(iv) $\beta$-D ribose sugar is present in nucleotide of DNA.\\
(v) Phosphodiester linkage is present in DNA.\\
(vi) DNA fragment\\ \qfigi{p1_chem_s3_q06_1}\\ .\\
(vii) DNA nucleic acid evolved before RNA nucleic acid in nature.\\
(viii) deoxy Adenosine, deoxy guanosine, deoxy thymidine and deoxy cytidine nucleoside are present in nucleotides of DNA.\\
(ix) Many genes are present in DNA which have information of specific protein in form of specific sequence of nucleotides.

\qn{7} Hydrolysis of tetrapeptide gives two asparagine, and two glycine amino acids. Calculate the value of X.\\
If \ X = 0.3 Z\\
Where Z is \% of nitrogen by weight in tetrapeptide.

\qn{8}
\qfig{p1_chem_s3_q08_1}
How many total 2$^\circ$ and 3$^\circ$ alcohol with molecular formula C$_{15}$H$_{16}$O can perform the above reactions.

\partsec{3}{MATHEMATICS}
\sect{Section-I(i)}

\qn{1} If $\displaystyle\int\dfrac{\left(\frac1x+\frac1{x^2}\right)(x-1)\,dx}{\left(\frac1{x^4}+\frac1{x^2}\right)\sqrt{(x^4-x^3+x^2)(x^4+x^3+x^2)}}=\sec^{-1}(f(x))+c$, $f(1)=2$ then local minimum value of f(x) is\\
(Where `c' is the constant of integration)
\begin{qoptions}
\opt{A}{2}
\opt{B}{5}
\opt{C}{6}
\opt{D}{8}
\end{qoptions}

\qn{2} Consider three sets $E_1=\{1,3,4\}$, $F_1=\{1,4,5\}$ and $G_1=\{3,4,5,6\}$. Two elements are chosen at random, without replacement, from the set $E_1$, and let $S_1$ denote the set of these chosen elements.\\
Let $E_2=E_1-S_1$ and $F_2=F_1\cup S_1$. Now two elements are chosen at random, without replacement, from the set $F_2$ and let $S_2$ denote the set of these chosen elements.\\
Let $G_2=G_1\cup S_2$. Finally, two elements are chosen at random, without replacement, from the set $G_2$ and let $S_3$ denote the set of these chosen elements.\\
Let $E_3=E_2\cup S_3$. Given that $E_1=E_3$, let p be the conditional probability of the event $S_1=\{1,3\}$. Then the value of p is
\begin{qoptions}
\opt{A}{$\dfrac45$}
\opt{B}{$\dfrac35$}
\opt{C}{$\dfrac15$}
\opt{D}{$\dfrac25$}
\end{qoptions}

\qn{3} Let $\theta_1,\theta_2,\dots,\theta_{100}$ be positive valued angles (in radian) such that $\displaystyle\sum_{i=1}^{100}\theta_i=2\pi$. Define the complex numbers $z_1=e^{i\theta_1}$ and $z_k|z_{k-1}|=2^{k-1}z_{k-1}e^{i\theta_k}$ for k = 2, 3, ..., 100, where $i=\sqrt{-1}$. Consider the statements P and Q given below :\\
P : $|z_2-z_1|+|z_3-z_2|+\dots+|z_{100}-z_{99}|+|z_1-z_{100}|\ge2^{101}-2$\\
Q : $\left|2^{98}z_2-2^{99}z_1\right|+\left|2^{97}z_3-2^{98}z_2\right|+\dots\left|z_{100}-2z_{99}\right|+\left|2^{99}z_1-z_{100}\right|\geqslant2^{100}\pi$
\begin{qoptions}
\opt{A}{P is \textbf{TRUE} and Q is \textbf{FALSE}}
\opt{B}{Q is \textbf{TRUE} and P is \textbf{FALSE}}
\opt{C}{both P and Q are \textbf{TRUE}}
\opt{D}{both P and Q are \textbf{FALSE}}
\end{qoptions}

\qn{4} Let a, b and c are real numbers such that the system of linear equations\\
$x+y+z=a$\\
$x+5y-z=b$\\
$x-y+2z=c+2$ is consistent.\\
Let $|M|$ represent the determinant of the matrix $M=\begin{bmatrix}a&-2&b\\2&1&c\\1&0&3\end{bmatrix}$. Let `P' be the plane containing all those (a, b, c) for which above system of linear equations is consistent and `D' be the square of distance of the point (1, 1, 1) from the plane `P'. Then the value of $|M|+14D$ is
\begin{qoptions}
\opt{A}{32}
\opt{B}{16}
\opt{C}{20}
\opt{D}{12}
\end{qoptions}

\sect{Section-I(ii)}

\qn{5} Let $f:\left(0,\dfrac\pi2\right)\to R$ satisfies differential equation $xf'(x)+(x\tan x+1)f(x)=\sec x\ \&\ f\left(\dfrac\pi4\right)=0$, then choose correct option(s)
\begin{qoptions}
\opt{A}{$f\left(\dfrac\pi3\right)=\dfrac{3\left(\sqrt3-1\right)}{2\pi}$}
\opt{B}{$f(x)=1$ has a solution}
\opt{C}{Area bounded by $g(x)=xf(x)$, x-axis from $x=\dfrac\pi4$ to $x=\dfrac\pi2$ is $\left(\sqrt2-1\right)$ square units}
\opt{D}{$g(x)=xf(x)$ is a bounded function}
\end{qoptions}

\qn{6} Consider two circles $S_1:x^2+y^2-8x=0$ and $S_2:x^2+y^2+2x=0$ -
\begin{qoptions}
\opt{A}{Number of common tangents of $S_1$ and $S_2$ is 3}
\opt{B}{Radical axis of $S_1$ and $S_2$ is $x=3$}
\opt{C}{Equation of circumcircle of the triangle formed by point of intersection of direct common tangents and the points where common tangents touches $S_2$ is $x^2+y^2+\dfrac{11}3x+\dfrac83=0$}
\opt{D}{Equation of circumcircle of the triangle formed by point of intersection of direct common tangents and the points where common tangents touches $S_1$ is $x^2+y^2+\dfrac83x+\dfrac{11}3=0$}
\end{qoptions}

\qn{7} Let $f:R\to(0,1)$ be a continuous function. Then which of the following function(s) must has(have) the value zero at some point in the interval (0, 1) ?
\begin{qoptions}
\opt{A}{$x^{10}-f(2x)$}
\opt{B}{$x^2-\displaystyle\int_0^xf(t)\,dt$}
\opt{C}{$\left(\displaystyle\int_0^{\frac\pi2-x}f(3t)\sin t\,dt\right)-x$}
\opt{D}{$\sqrt{x^3}-\displaystyle\int_0^x\sqrt[3]{f(t)}\,dt$}
\end{qoptions}

\qn{8} Let $g(x)=\begin{cases}2(x+1)&,\ -\infty<x\leqslant-1\\ \sqrt{1-x^2}&,\ -1<x<1\\ \big|\big||x|-1\big|-1\big|&,\ 1\leqslant x<\infty\end{cases}$, then choose the correct option(s) :
\begin{qoptions}
\opt{A}{$g(x)$ is non-differentiable at exactly three points}
\opt{B}{$g(x)$ is continuous in $(-\infty,1]$}
\opt{C}{$g(x)$ is differentiable in $(-\infty,-1)$}
\opt{D}{$g(x)$ has discontinuity at $x=1$ but continuous at $x=-1$}
\end{qoptions}

\qn{9} If $z^{11}=1$, then which of the following is(are) true ?
\begin{qoptions}
\opt{A}{There are exactly 5 complex numbers for which $\cos(\arg(z))<0$}
\opt{B}{The value of $|1-\alpha_i|^2+|1+\alpha_i|^2=4$ where $\alpha_i$ is non-real roots of $z^{11}=1$ and i = 1, ....10}
\opt{C}{The range of $|1+z|$ contains 6 distinct elements}
\opt{D}{For non$-$real z, $\arg\left(\dfrac{1+z}{1-z}\right)$ contains 10 distinct elements}
\end{qoptions}

\qn{10} Let A and B be two square matrices of order 3 then which of the following statement(s) is(are) correct ?
\begin{qoptions}
\opt{A}{$ABA^T$ is symmetric matrix}
\opt{B}{$AB-BA$ is skew symmetric matrix}
\opt{C}{If $B=|A|\,A^{-1}$, $|A|\neq0$ then $\text{adj.}(A^T)-B$ is a skew symmetric matrix}
\opt{D}{If $B+A^T=\mathbf O$ and A is skew symmetric matrix then $B^{15}$ is also skew symmetric matrix}
\end{qoptions}

\sect{Section-III}

\qn{1} If N denotes the number of positive integral solutions of equation $xy=2^3.3^4.5^6.(x+y)$, then $\dfrac N{91}$ is equal to

\qn{2} Three distinct numbers are chosen one by one without replacement form the set $\{1,2,3,\dots\dots,9\}$. If it is known that sum of first two chosen numbers is even, then $\dfrac pq$ is the probability that sum of all three numbers is even. Then $(q-2p)$ is equal to (where p \& q are coprime numbers)

\qn{3} Consider the parabola $y^2=4x$. Let P and Q be two points $(4,-4)$ and $(9,6)$ on the parabola. Let R be a moving point on the arc of the parabola between P and Q. If `A' denote the maximum area of $\Delta RPQ$ then $\left[\dfrac{\sqrt5A}{2}\right]$ is equal to (where [.] denotes greatest integer function)

\qn{4} Evaluate\\
\resizebox{\linewidth}{!}{$\displaystyle\lim_{x\to0}\ \lim_{n\to\infty}\dfrac{\left[\begin{array}{@{}c@{}}(n+x)(n+2x)\dots(n+2nx)\\ \times\left(\tan^2\tfrac{\pi}{2n}.\tan^2\tfrac{2\pi}{2n}.\tan^2\tfrac{3\pi}{2n}\dots\tan^2\tfrac{(n-1)\pi}{2n}\right)\end{array}\right]^{\frac1{2n}}}{n-\sqrt n}$}\\

\qn{5} Let $S_1$ and $S_2$ are focii of the hyperbola with length of transverse axis equal to 20. Perpendicular distance of $S_1$ from normal at P is 20. If the normal at P divides the line segment joining $S_1$ and $S_2$ in ratio 5 : 1 externally. If the square of length of semi-conjugate axis is $\lambda$, then the value of $\dfrac\lambda5$ is

\qn{6} Suppose $\det\begin{bmatrix}\displaystyle\sum_{k=1}^{n}k&\displaystyle\sum_{k=1}^{n}{}^nC_kk^2\\ \displaystyle\sum_{k=1}^{n}{}^nC_kk&\displaystyle\sum_{k=0}^{n}{}^nC_k3^k\end{bmatrix}+\dfrac{n(n+1).4^n}2=0$, holds for some positive integer n. Then $\dfrac9{251}\displaystyle\sum_{k=1}^{n}\dfrac{{}^nC_k}{k+1}$ equals

\qn{7} Let $\vec a,\vec b$ and $\vec c$ be vectors in three-dimensional space, where $\vec a.\vec a=1$, $\vec b.\vec b=2$, $\vec c.\vec c=9$, $\vec a.\vec c=2$, $\vec b.\vec c=1$. If the volume of the parallelopiped, whose adjacent sides are represented by the vectors $\vec a,\vec b$ and $\vec c$, is $\dfrac{\sqrt{85}}3$, then the value of $\dfrac9{83}\left|\vec a+\vec b+\vec c\right|^2$ is

\qn{8} In a class of 5 students average weight of 4 lightest students is 40 kgs, Average weight of 4 heaviest students is 45 kgs. The difference between maximum and minimum possible average weight overall (in kg) is

\end{multicols}
\end{document}
